\begin{exo}[CNS valeur propre commune]
	\label{reduction1}
	Soient $A,B\in\mathcal M_n(\C)$. Montrer que les propositions suivantes sont équivalentes :
	\begin{enumerate}
		\item $A$ et $B$ possèdent une valeur propre commune.
		\item Il existe $M\in\mathcal M_n(\C)$ non nulle telle que $AM=MB$
		\item $\mu_A(B)\notin \gl_n(\C)$
	\end{enumerate}
\end{exo}

\begin{solution}
    \begin{itemize}
        \item[]$\boxed{\text{1.} \Rightarrow \text{2.}}$\quad
        Supposons que $A$ et $B$ possèdent une valeur propre commune. Notons $\lambda$ cette valeur propre.
        D'abord, il existe $X\in\mathcal M_{n,1}(\C)$ une colonne non nulle telle que
        \[AX=\lambda X\]
        Mais $\lambda$ étant valeur propre de $B$, elle est encore valeur propre de ${}^tB$, d'où l'existence de $Y\in\mathcal M_{n,1}(\C)$ une colonne non nulle telle que
        \[{}^tBY=\lambda Y\]
        On pose $M=X{}^tY$. On vérifie facilement que $AM=MB$, puis $M$ est non nulle puisque de rang $1$.
        
        \item[]$\boxed{\text{2.} \Rightarrow \text{3.}}$\quad
        Déjà, on a 
        \[A^2M=A(AM)=(AM)B=MB^2\]
        puis, on vérifie facilement par récurrence que 
        \[\forall k\in\N,\ A^kM=MB^k\]
        D'où, pour tout polynôme $P\in\C[X]$, l'égalité
        \[P(A)M=MP(B)\]
        En particulier, on a 
        \[M\mu_A(B)=\mu_A(A)M=0\]
        Ainsi, $M$ étant non nulle, $\mu_A(B)$ est soit nulle, soit un diviseur de $0$, donc est non inversible.
        
        \item[]$\boxed{\text{3.} \Rightarrow \text{1.}}$\quad
        Par contraposée. Supposons que $A$ et $B$ n'ont pas de valeurs propres communes.
        Notons $\Sp(A)=\lbrace \lambda_1,\dots,\lambda_p\rbrace$ le spectre de $A$. On écrit
        \[\mu_A=(X-\lambda_1)^{\alpha_1}\dots(X-\lambda_p)^{\alpha_p}\]
        Or, pour tout $i\in\inl1p$, $\lambda_i$ n'est pas valeur propre de $B$ et donc $B-\lambda_iI_n\in \gl_n(\C)$. Ainsi $\mu_A(B)$ est un produit de matrices inversibles et est donc inversible.
    \end{itemize}
\end{solution}

\begin{exo}[$P(A)$ diagonalisable et $P'(A)$ inversible $\implies$ $A$ diagonalisable]
	\label{reduction2}
    Soit $A\in\mathcal M_n(\C)$ et $P\in\C[X]$ tel que $P(A)$ est diagonalisable et $P'(A)$ est inversible. Montrer que $A$ est diagonisable.
\end{exo}

\begin{solution}
    On note $B=P(A)$ et $\mu_B$ le polynôme minimal de $B$. 
    Le polynôme $\mu_B\circ P$ annule $A$, donc $\mu_A$ divise $\mu_B\circ P$. Il existe ainsi $Q\in\C[X]$ tel que $\mu_B\circ P = Q\mu_A$.
    Dérivons cette égalité polynomiale :
    \[(\mu_B'\circ P)P' = Q\mu_A' + Q'\mu_A \tag{*}\]
    Soit $\lambda \in \Sp(A)$ une valeur propre de $A$. Comme $\lambda$ est racine de $\mu_A$, l'évaluation de l'égalité (*) en $\lambda$ donne :
    \[\mu_B'(P(\lambda))P'(\lambda) = Q(\lambda)\mu_A'(\lambda)\]
    Puisque $B = P(A)$ est diagonalisable, son polynôme minimal $\mu_B$ est scindé à racines simples (SRS). Or, si $\lambda$ est valeur propre de $A$, $P(\lambda)$ est valeur propre de $P(A) = B$, et est donc une racine de $\mu_B$. La racine étant simple, on a $\mu_B'(P(\lambda)) \neq 0$.
    Par ailleurs, $P'(A)$ étant inversible, $0$ n'est pas valeur propre de $P'(A)$. Les valeurs propres de $P'(A)$ étant les $P'(\lambda_i)$, on en déduit que $P'(\lambda) \neq 0$.
    Ainsi, le membre de gauche de l'égalité est non nul. On en déduit que :
    \[Q(\lambda)\mu_A'(\lambda) \neq 0 \implies \mu_A'(\lambda) \neq 0\]
    Cette conclusion étant valide pour toute racine $\lambda$ de $\mu_A$, le polynôme $\mu_A$ est à racines simples. $A$ est donc diagonalisable.
\end{solution}

\begin{exo}[Diagonalisation dans $\mathcal M_n(\mathbb F_p)$]
	\label{reduction3}
    Soit $p$ un nombre premier, $n\in\N^*$ et $A\in\mathcal M_n(\mathbb F_p)$. Montrer que \[A\quad \text{diagonalisable}\quad \iff\quad A^p=A\]
\end{exo}

\begin{solution}
\begin{itemize}
\item[$\boxed{\Rightarrow}$] Supposons $A$ diagonalisable. Il existe alors $P\in \gl_n(\mathbb F_p)$ et $D\in \mathcal D_n(\mathbb F_p)$ telles que
\[A=PDP^{-1}\]
Mais alors $A^p-A = PD^pP^{-1} - PDP^{-1}$. Or, $D$ étant diagonale, on peut noter : 
\[ D = \begin{pmatrix} \lambda_1 & 0 & 0 \\ 0 & \ddots & 0 \\ 0 & 0 & \lambda_n \end{pmatrix} \]
On a immédiatement :
\[ D^p = \begin{pmatrix} \lambda_1^p & 0 & 0 \\ 0 & \ddots & 0 \\ 0 & 0 & \lambda_n^p \end{pmatrix} \]
D'après le petit théorème de Fermat (ou le théorème de Lagrange sur $\mathbb F_p^*$), pour tout $\lambda\in\mathbb F_p$, on a $\lambda^p=\lambda$. Donc $D^p=D$, d'où $A^p=A$.
\item[$\boxed{\Leftarrow}$] Supposons $A^p=A$. Le polynôme $P = X^p-X \in \mathbb F_p[X]$ est donc un polynôme annulateur de $A$. 
On factorise $P = X(X^{p-1}-1)$. En notant $Q = X^{p-1}-1$, on a d'après Fermat :
\[\forall x\in\mathbb F_p^*,\quad Q(x)=0\]
Le polynôme $Q$, de degré $p-1$, admet donc $p-1$ racines distinctes (tous les éléments de $\mathbb F_p^*$). Comme il est unitaire, il se factorise exactement en $Q = \prod_{x\in\mathbb F_p^*}(X-x)$.
On en déduit que :
\[P = X \prod_{x\in\mathbb F_p^*}(X-x) = \prod_{x\in\mathbb F_p}(X-x)\]
Le polynôme annulateur $P$ est scindé et à racines simples sur $\mathbb F_p$. Ainsi, $A$ est diagonalisable.
\end{itemize}
\end{solution}

\begin{exo}[Matrice semblable à son double]
	\label{reduction4}
	Soit $M\in\mathcal M_n(\C)$ telle que $M\underset{\text{sb}}{\sim}2M$. Montrer que $M$ est nilpotente.
\end{exo}

\begin{solution}
	On trigonalise $M$. Il existe $P\in \gl_n(\C)$ et $T\in \mathcal T_n(\C)$ tels que $M=PTP^{-1}$. 
    Puis, comme $M$ est semblable à $2M$, elle-même semblable à $2T$ (il suffit de multiplier par $2$ plus haut), $T$ est semblable à $2T$. 
    Donc $T$ et $2T$ ont même polynôme caractéristique. On note $\Sp(T)=\lbrace \lambda_1,\dots,\lambda_p\rbrace$ le spectre de $T$, $\chi_T$ 
    le polynôme caractéristique de $T$. Soit $\lambda\in \text{Sp}(T)$, alors $\chi_T(2\lambda)=\chi_{2T}(2\lambda)=0$, donc $2\lambda \in \text{Sp}(T)$. 
    Par récurrence sur $n\in\N$, on a \[\forall \lambda \in \text{Sp}(T),\ \forall n\in\N,\ 2^n\lambda\in\text{Sp}(T)\]
	Ceci est exclu lorsqu'il existe une valeur propre non nulle de $T$. 
    Ainsi, toutes les valeurs propres sont nulles, $T$ est de diagonale nulle (et triangulaire) et est donc nilpotente, puis $M$ lui étant semblable, 
    elle est également nilpotente.
\end{solution}

\begin{exo}[Comparaison de polynômes minimaux]
	\label{reduction5}
    Soit $E$ un $K$ espace vectoriel, $f\in L(E)$ et $G:g\in L(E)\mapsto f\circ g$. Vérifier que $G\in L(L(E))$ et comparer (sous réserve d'existence) les polynômes minimaux de $f$ et $G$.
\end{exo}

\begin{solution}
	On laisse au lecteur les soins de vérifier que $G$ est bien un endomorphisme. 
    Supposons que $\mu_f$ existe. 
    Soit $P\in K[X]$. 
    Remarquons que pour tout $g\in L(E)$, $P(G)(g)=P(f)\circ g$ (on peut d'abord montrer cela pour $P=X^k$ pour tout $k$ entier naturel par récurrence, 
    et on passe ensuite à tout polynôme par des combinaisons linéaires). Ainsi, $P(f)=0\implies P(G)=0$. 
    Ceci montre deux choses : d'abord, $G$ admet un polynôme annulateur non nul, donc $\mu_G$ existe, et $(\mu_f)\subset (\mu_G)$, soit $\mu_G$ divise $\mu_f$. 
\end{solution}

\textit{Les deux exercices qui vont suivre n'ont jamais été posés lors de mon année scolaire, mais ont fait partie de mes réflexions lors de mon année de MPI*, j'ai alors tenu à les inclure lors de la première rédaction de ce document.}

\begin{exo}[Coefficients du polynôme caractéristique]
	\label{reduction6}
	Soit $M\in\mathcal M_n(\K)$. On appelle mineur principal d'ordre $k\in\lbrace1,\dots,n\rbrace$ le déterminant d'un $M_I=(m_{i,j})_{(i,j)\in I^2}$ avec $I\subset\lbrace 1,\dots,n\rbrace$ tel que $\Card(I)=k$. Donner une expression des coefficients de $\chi_M$ en fonction des mineurs principaux.
\end{exo}

\begin{solution}
    Correction trouvable sur \href{https://perso.eleves.ens-rennes.fr/people/amar.ahmane}{ma page personnelle}.
\end{solution}

\begin{exo}[Limite d'une suite de matrices]
	\label{reduction7}
	Soit $M\in\mathcal M_n(\K)$, $\K = \R$ ou $\C$. Donner une condition nécessaire et suffisante pour que la suite $(M^p)_{p\in\N}$ converge. Que dire sur la valeur de la limite en cas de convergence ?
\end{exo}

\begin{solution}
    On procède par analyse synthèse :
    
    \textbf{Analyse :} Supposons que la suite $(M^p)$ converge, notons $L$ sa limite.
    Remarquons qu'en écrivant $(M^p)^2=M^{2p}$ pour tout $p$ et en passant à la limite, on obtient $L^2=L$,
    $L$ est alors la matrice d'une projection. Remarquons de plus qu'en écrivant $MM^p=M^pM$ pour tout $p$ et en passant à la limite, on obtient $ML=LM$.
    Ainsi, comme $M$ et $L$ commutent, il existe $P\in\gl_n(\C)$ et $T_L$ et $T_M$ des matrices trigonales supérieures de $\mathcal M_n(\C)$ telles que 
    $M=PT_MP^{-1}$ et $L=PT_LP^{-1}$. 
    Comme $M^p\to L$, on a en fait $T_M^p\to T_L$. 
    Ainsi pour tout $i,j=1,\dots,n$, $[T_M^p]_{i,j}\to [T_L]_{i,j}$.
    En regardant les coefficients sur la diagonale, on a pour tout $i=1,\dots,n$, $\lambda_i^p\to [T_L]_{i,i}$
    où $\lambda_1,\dots,\lambda_n$ sont les valeurs propres complexes de $M$ dans l'ordre dans lequel elles apparaissent dans $T_M$.
    $L$ étant une matrice de projection, ses valeurs propres sont soit $1$ soit $0$, donc pour tout $i=1,\dots,n$,
    $\lambda_i^p\to 0$ ou $\lambda_i^p\to 1$.
    Soit $i\in\lbrace1,\dots,n\rbrace$.
    \begin{itemize}
        \item Si $\lambda_i^p\to 0$, alors $|\lambda_i|<1$ trivialement ;
        \item Si $\lambda_i^p\to 1$, alors $|\lambda_i|=1$. 
        Montrons qu'en fait $\lambda_i=1$. 
        Soit $\theta\in\R$ tel que $e^{i\theta}=\lambda_i$.
        On a $|2\sin(p\theta/2)|=|e^{ip\theta}-1|\to 0$. 
        On procède par l'absurde et on suppose que $\theta \notin 2\pi\Z$.
        Deux cas de figure se présentent, soit $\theta$ est commensurable à $\pi$ 
        (i.e on peut écrire $\theta=k\pi/q$ avec $k,q$ des entiers, $q>0$),
        et alors, comme il faut $\theta\notin2\pi\Z$, on a $\theta/2\notin\pi\Z$, 
        et alors pour tout $p$, 
        $\sin((4qp+1)\theta/2)=\sin(2kp\pi+\theta/2)=\sin(\theta/2)\neq 0$,
        et on exhibe ainsi une sous-suite de $(\sin(p\theta/2))$ qui ne tend pas vers $0$ ;
        soit $\theta$ n'est pas commensurable à $\pi$, et alors le sous-groupe $(\theta/2)\Z+2\pi\Z$ de $(\R,+)$ est dense dans $\R$.
        Il vient alors, par continuité de $\sin$, que $\sin((\theta/2)\Z)=\sin((\theta/2)\Z+2\pi\Z)$ est dense dans $\sin(\R)=[-1,1]$, et on trouve alors 
        que l'ensemble des termes de la suite $(|\sin(p\theta/2)|)_p$ est dense dans $[0,1]$, alors que cette suite converge... 
        On obtient alors que $\theta\in 2\pi\Z$, soit que $\lambda_i=1$.
    \end{itemize}
    
    \textbf{Synthèse :} Soit à présent une matrice $M$ dont toutes les valeurs propres complexes sont soit égales à $1$, soit de module strictement plus petit que $1$.
    La décomposition de Dunford donne $D$ diagonalisable et $N$ nilpotente telles que $M=N+D$ et $ND=DN$.
    Pour tout $p\geq n$, on a 
    \[
        M^p=(N+D)^p=\sum_{k=0}^p\binom pkD^{p-k}N^k=\sum_{k=0}^{n-1}\binom pk D^{p-k}N^k   \tag*{(*)} 
    \]
    Soit $P\in\gl_n(\C)$ telle que $\Delta=PDP^{-1}$ soit diagonale. 
    Dans un premier temps, supposons que $1$ n'est pas valeur propre de $M$.
    Dans ce cas, pour tout $k\in\lbrace0,\dots,n-1\rbrace$, $\binom pk\Delta^{p-k}\to 0$ : en effet,
    ces matrices sont diagonales, donc pour $i=1,\dots,n$, $\left[\binom pk\Delta^{p-k}\right]_{i,i}=\binom pk\lambda_i^{p-k}$ avec 
    $|\lambda_i|<1$. On a alors $\left[\binom pk\Delta^{p-k}\right]_{i,i}\to 0$, car $\binom pk\sim p^k/k!$ et, par croissances comparées, $p^k=o\left(1/|\lambda_i|^{p-k}\right)$.
    Par continuité de la multiplication par une matrice $\binom pk D^{p-k}=\binom pkP\Delta^{p-k}P^{-1}\to 0$ et alors $M^p\to 0$ comme somme finie de termes qui tendent vers $0$.
    
    \textit{À ce stade, tout peut plus ou moins être fait dans le cadre du programme (ou en débordant un tout petit peu, l'existence de la décomposition de Dunford n'est pas difficile à établir). Il reste à traiter le cas où $1$ est valeur propre, et ce cas est plus difficile et demandera d'utiliser un résultat de réduction qui est hors programme, celui de la réduction de Jordan.}
   
   
    Lorsque $1$ est valeur propre de $M$, on voit que dès que la multiplicité de $1$ dans le polynôme minimal est plus grande que $2$, alors il existe un bloc de Jordan lié à $1$ de taille au moins $2$. On en déduit que la suite des puissances de $M$ ne converge pas. On peut également montrer que $(M^p)$ converge lorsque $1$ est racine simple du polynôme minimal de $M$.
\end{solution}